% !TEX program = xelatex
\documentclass[12pt]{article}
\usepackage[paperwidth=612bp,paperheight=792bp,margin=0bp,noheadfoot]{geometry}
\usepackage{fontspec}
\usepackage{xeCJK}
\usepackage{tikz}
\pagestyle{empty}
\setlength{\parindent}{0pt}
\setlength{\parskip}{0pt}
\IfFontExistsTF{TeX Gyre Termes}{\setmainfont{TeX Gyre Termes}}{\setmainfont{Times New Roman}}
\IfFontExistsTF{Noto Serif CJK SC}{\setCJKmainfont{Noto Serif CJK SC}}{\setCJKmainfont{Songti SC}}
\newcommand{\QuestionTitle}{人机实验1 H-C3：高祖本纪}
\newcommand{\DrawQuestionTitle}{%
\node[anchor=base west,inner sep=0pt,outer sep=0pt] at ([xshift=72bp,yshift=-34bp]current page.north west) {{\fontsize{14bp}{14bp}\selectfont\bfseries \QuestionTitle}};%
\node[anchor=base west,inner sep=0pt,outer sep=0pt] at ([xshift=72bp,yshift=-62bp]current page.north west) {{\fontsize{12bp}{12bp}\selectfont 用时：\underline{\hspace{80bp}}}};%
}

\begin{document}
% Auto-generated for 刘邦
% Keep only 8.5.1--8.6.1 and 8.22.1--8.26.3.
\thispagestyle{empty}
\begin{tikzpicture}[remember picture,overlay]
\DrawQuestionTitle
\node[anchor=base west,inner sep=0pt,outer sep=0pt] at ([xshift=72bp,yshift=-84bp]current page.north west) {{\fontsize{12bp}{12bp}\selectfont 【8.5.1】 高祖作亭长时，曾经请假回家。}};
\node[anchor=base west,inner sep=0pt,outer sep=0pt] at ([xshift=72bp,yshift=-99bp]current page.north west) {{\fontsize{12bp}{12bp}\selectfont 【8.5.2】 吕后与两个孩子在田间除草，有一老人路过，要些水喝，吕后就请他吃了饭。}};
\node[anchor=base west,inner sep=0pt,outer sep=0pt] at ([xshift=72bp,yshift=-114bp]current page.north west) {{\fontsize{12bp}{12bp}\selectfont 【8.5.3】 老人给吕后相面，说：“夫人是天下的贵人。”吕后让他给两个孩子看相。}};
\node[anchor=base west,inner sep=0pt,outer sep=0pt] at ([xshift=72bp,yshift=-129bp]current page.north west) {{\fontsize{12bp}{12bp}\selectfont 【8.5.4】 老人看了孝惠，说：“夫人所以显贵，就是这个孩子的缘故。”看了鲁元，也}};
\node[anchor=base west,inner sep=0pt,outer sep=0pt] at ([xshift=72bp,yshift=-144bp]current page.north west) {{\fontsize{12bp}{12bp}\selectfont 是贵相。}};
\node[anchor=base west,inner sep=0pt,outer sep=0pt] at ([xshift=72bp,yshift=-159bp]current page.north west) {{\fontsize{12bp}{12bp}\selectfont 【8.5.5】 老人已经走了，高祖正好从别人家来到田间，吕后告诉他一位客人从这里经过}};
\node[anchor=base west,inner sep=0pt,outer sep=0pt] at ([xshift=72bp,yshift=-174bp]current page.north west) {{\fontsize{12bp}{12bp}\selectfont ，给我们母子看相，说将来都是大贵人，高祖问老父在哪儿.}};
\node[anchor=base west,inner sep=0pt,outer sep=0pt] at ([xshift=72bp,yshift=-189bp]current page.north west) {{\fontsize{12bp}{12bp}\selectfont 【8.5.6】 吕后说：“走出不远。”}};
\node[anchor=base west,inner sep=0pt,outer sep=0pt] at ([xshift=72bp,yshift=-204bp]current page.north west) {{\fontsize{12bp}{12bp}\selectfont 【8.5.7】 高祖追上了老人，向他询问。}};
\node[anchor=base west,inner sep=0pt,outer sep=0pt] at ([xshift=72bp,yshift=-219bp]current page.north west) {{\fontsize{12bp}{12bp}\selectfont 【8.5.8】 老人说：“刚才相过夫人和孩子，他们都跟你相似，你的相貌，贵不可言。”}};
\node[anchor=base west,inner sep=0pt,outer sep=0pt] at ([xshift=72bp,yshift=-234bp]current page.north west) {{\fontsize{12bp}{12bp}\selectfont 【8.5.9】 高祖便道谢说：“如果真像老父所说，决不忘记对我的恩德。”}};
\node[anchor=base west,inner sep=0pt,outer sep=0pt] at ([xshift=72bp,yshift=-249bp]current page.north west) {{\fontsize{12bp}{12bp}\selectfont 【8.5.10】 等到高祖显贵，竟然不知道老人的去处了。}};
\node[anchor=base west,inner sep=0pt,outer sep=0pt] at ([xshift=72bp,yshift=-264bp]current page.north west) {{\fontsize{12bp}{12bp}\selectfont 【8.6.1】 高祖做亭长，以竹皮为帽，这帽子是他派求盗到薛县制做的，经常戴着它。}};
\node[anchor=base west,inner sep=0pt,outer sep=0pt] at ([xshift=72bp,yshift=-279bp]current page.north west) {{\fontsize{12bp}{12bp}\selectfont 【8.22.1】 汉元年十月，沛公的军队先于各路诸侯到达霸上。}};
\node[anchor=base west,inner sep=0pt,outer sep=0pt] at ([xshift=72bp,yshift=-294bp]current page.north west) {{\fontsize{12bp}{12bp}\selectfont 【8.22.2】 秦王于婴素车白马，用丝带系着脖子，封了皇帝的印玺和符节，在轵道旁投降}};
\node[anchor=base west,inner sep=0pt,outer sep=0pt] at ([xshift=72bp,yshift=-309bp]current page.north west) {{\fontsize{12bp}{12bp}\selectfont 。}};
\node[anchor=base west,inner sep=0pt,outer sep=0pt] at ([xshift=72bp,yshift=-324bp]current page.north west) {{\fontsize{12bp}{12bp}\selectfont 【8.22.3】 将领们有的主张杀死秦王。}};
\node[anchor=base west,inner sep=0pt,outer sep=0pt] at ([xshift=72bp,yshift=-339bp]current page.north west) {{\fontsize{12bp}{12bp}\selectfont 【8.22.4】 沛公说：“当初楚怀王派遣我，本来是因为我能宽大容人。况且人家已经降服}};
\node[anchor=base west,inner sep=0pt,outer sep=0pt] at ([xshift=72bp,yshift=-354bp]current page.north west) {{\fontsize{12bp}{12bp}\selectfont ，又杀死人家，不吉利。”}};
\node[anchor=base west,inner sep=0pt,outer sep=0pt] at ([xshift=72bp,yshift=-369bp]current page.north west) {{\fontsize{12bp}{12bp}\selectfont 【8.22.5】 于是就把秦王交给了官吏，向西进入咸阳。}};
\node[anchor=base west,inner sep=0pt,outer sep=0pt] at ([xshift=72bp,yshift=-384bp]current page.north west) {{\fontsize{12bp}{12bp}\selectfont 【8.22.6】 沛公想要留在宫殿中休息，樊哙、张良劝说后，才封闭了秦宫的贵重珍宝、财}};
\node[anchor=base west,inner sep=0pt,outer sep=0pt] at ([xshift=72bp,yshift=-399bp]current page.north west) {{\fontsize{12bp}{12bp}\selectfont 物和库房，回军霸上。}};
\node[anchor=base west,inner sep=0pt,outer sep=0pt] at ([xshift=72bp,yshift=-414bp]current page.north west) {{\fontsize{12bp}{12bp}\selectfont 【8.22.7】 召集各县的父老、豪杰说：“父老们苦于秦朝的严刑峻法已经很久了，诽谤朝}};
\node[anchor=base west,inner sep=0pt,outer sep=0pt] at ([xshift=72bp,yshift=-429bp]current page.north west) {{\fontsize{12bp}{12bp}\selectfont 政的要灭族，相聚议论的要在街市上处斩。我和诸侯们约定，先入关的在关中称王，我应}};
\node[anchor=base west,inner sep=0pt,outer sep=0pt] at ([xshift=72bp,yshift=-444bp]current page.north west) {{\fontsize{12bp}{12bp}\selectfont 当称王关中。同父老们约定，法律只有三章：杀人的处死，伤人和抢劫的处以与所犯罪相}};
\node[anchor=base west,inner sep=0pt,outer sep=0pt] at ([xshift=72bp,yshift=-459bp]current page.north west) {{\fontsize{12bp}{12bp}\selectfont 当的刑罚。其余的秦朝法律全部废除。官吏和百姓都要安居如故。我所以到这里来，是为}};
\node[anchor=base west,inner sep=0pt,outer sep=0pt] at ([xshift=72bp,yshift=-474bp]current page.north west) {{\fontsize{12bp}{12bp}\selectfont 父老们除害，不会有欺凌暴虐的行为，不要害怕。我所以回军霸上，是等待诸侯们到来制}};
\node[anchor=base west,inner sep=0pt,outer sep=0pt] at ([xshift=72bp,yshift=-489bp]current page.north west) {{\fontsize{12bp}{12bp}\selectfont 定共同遵守的纪律。”}};
\node[anchor=base west,inner sep=0pt,outer sep=0pt] at ([xshift=72bp,yshift=-504bp]current page.north west) {{\fontsize{12bp}{12bp}\selectfont 【8.22.8】 沛公派人与秦朝官吏巡行县城乡间，告谕百姓。}};
\node[anchor=base west,inner sep=0pt,outer sep=0pt] at ([xshift=72bp,yshift=-519bp]current page.north west) {{\fontsize{12bp}{12bp}\selectfont 【8.22.9】 秦地的百姓大为高兴，争先恐后地拿出牛羊酒食款待士兵。}};
\node[anchor=base west,inner sep=0pt,outer sep=0pt] at ([xshift=72bp,yshift=-534bp]current page.north west) {{\fontsize{12bp}{12bp}\selectfont 【8.22.10】 沛公又谦让不肯接受，说：“仓库的谷子很多，不缺乏，不愿破费百姓。”}};
\node[anchor=base west,inner sep=0pt,outer sep=0pt] at ([xshift=72bp,yshift=-549bp]current page.north west) {{\fontsize{12bp}{12bp}\selectfont 百姓更加高兴，唯恐沛公不做秦王。}};
\node[anchor=base west,inner sep=0pt,outer sep=0pt] at ([xshift=72bp,yshift=-564bp]current page.north west) {{\fontsize{12bp}{12bp}\selectfont 【8.23.1】 有人劝沛公说：“秦地比天下富足十倍，地势好。如今听说章邯投降了项羽，}};
\node[anchor=base west,inner sep=0pt,outer sep=0pt] at ([xshift=72bp,yshift=-579bp]current page.north west) {{\fontsize{12bp}{12bp}\selectfont 项羽就给了雍王的封号，称王于关中。现在即将来到关中就国，你沛公恐怕不能占有这个}};
\node[anchor=base west,inner sep=0pt,outer sep=0pt] at ([xshift=72bp,yshift=-594bp]current page.north west) {{\fontsize{12bp}{12bp}\selectfont 地方了。应赶快派兵把守函谷关，不让诸侯军进来，逐渐征集关中兵，以加强实力，抵抗}};
\node[anchor=base west,inner sep=0pt,outer sep=0pt] at ([xshift=72bp,yshift=-609bp]current page.north west) {{\fontsize{12bp}{12bp}\selectfont 诸侯兵。”}};
\node[anchor=base west,inner sep=0pt,outer sep=0pt] at ([xshift=72bp,yshift=-624bp]current page.north west) {{\fontsize{12bp}{12bp}\selectfont 【8.23.2】 沛公赞成他的计策，照着做了。}};
\node[anchor=base west,inner sep=0pt,outer sep=0pt] at ([xshift=72bp,yshift=-639bp]current page.north west) {{\fontsize{12bp}{12bp}\selectfont 【8.23.3】 十一月间，项羽果然率领诸侯军西进，想要入关，而关门闭着。}};
\node[anchor=base west,inner sep=0pt,outer sep=0pt] at ([xshift=72bp,yshift=-654bp]current page.north west) {{\fontsize{12bp}{12bp}\selectfont 【8.23.4】 听说沛公已经平定关中，大怒，派黥布等攻破了函谷关。}};
\node[anchor=base west,inner sep=0pt,outer sep=0pt] at ([xshift=72bp,yshift=-669bp]current page.north west) {{\fontsize{12bp}{12bp}\selectfont 【8.23.5】 十二月间，就到了戏水。}};
\node[anchor=base west,inner sep=0pt,outer sep=0pt] at ([xshift=72bp,yshift=-684bp]current page.north west) {{\fontsize{12bp}{12bp}\selectfont 【8.23.6】 沛公左司马曹无伤听说项王发怒，}};
\node[anchor=base west,inner sep=0pt,outer sep=0pt] at ([xshift=72bp,yshift=-699bp]current page.north west) {{\fontsize{12bp}{12bp}\selectfont 【8.23.7】 要攻打沛公，派人告诉项羽说：“沛公想要称王关中，令子婴为相，珍宝被他}};
\node[anchor=base west,inner sep=0pt,outer sep=0pt] at ([xshift=72bp,yshift=-714bp]current page.north west) {{\fontsize{12bp}{12bp}\selectfont 全部占有了。”}};
\end{tikzpicture}
\null
\newpage

\thispagestyle{empty}
\begin{tikzpicture}[remember picture,overlay]
\DrawQuestionTitle
\node[anchor=base west,inner sep=0pt,outer sep=0pt] at ([xshift=72bp,yshift=-84bp]current page.north west) {{\fontsize{12bp}{12bp}\selectfont 【8.23.8】 打算以此求得封赏。}};
\node[anchor=base west,inner sep=0pt,outer sep=0pt] at ([xshift=72bp,yshift=-99bp]current page.north west) {{\fontsize{12bp}{12bp}\selectfont 【8.23.9】 亚父劝项羽进攻沛公。}};
\node[anchor=base west,inner sep=0pt,outer sep=0pt] at ([xshift=72bp,yshift=-114bp]current page.north west) {{\fontsize{12bp}{12bp}\selectfont 【8.23.10】 当时项羽饱餐士卒，准备明日会战。}};
\node[anchor=base west,inner sep=0pt,outer sep=0pt] at ([xshift=72bp,yshift=-129bp]current page.north west) {{\fontsize{12bp}{12bp}\selectfont 【8.23.11】 这时项羽兵四十万，号称百万。}};
\node[anchor=base west,inner sep=0pt,outer sep=0pt] at ([xshift=72bp,yshift=-144bp]current page.north west) {{\fontsize{12bp}{12bp}\selectfont 【8.23.12】 沛公兵十万，号称二十万，兵力敌不过项羽。}};
\node[anchor=base west,inner sep=0pt,outer sep=0pt] at ([xshift=72bp,yshift=-159bp]current page.north west) {{\fontsize{12bp}{12bp}\selectfont 【8.23.13】 恰巧项伯要救张良，夜间去见他。}};
\node[anchor=base west,inner sep=0pt,outer sep=0pt] at ([xshift=72bp,yshift=-174bp]current page.north west) {{\fontsize{12bp}{12bp}\selectfont 【8.23.14】 （回来后，）用道理劝说项羽，项羽取消了进攻沛公的计划。}};
\node[anchor=base west,inner sep=0pt,outer sep=0pt] at ([xshift=72bp,yshift=-189bp]current page.north west) {{\fontsize{12bp}{12bp}\selectfont 【8.23.15】 沛公带来了一百多骑兵，驰至鸿门，来见项羽，表示歉意。}};
\node[anchor=base west,inner sep=0pt,outer sep=0pt] at ([xshift=72bp,yshift=-204bp]current page.north west) {{\fontsize{12bp}{12bp}\selectfont 【8.23.16】 项羽说：“这是你沛公左司马曹无伤向我说的。不然，我项羽何至于做这样}};
\node[anchor=base west,inner sep=0pt,outer sep=0pt] at ([xshift=72bp,yshift=-219bp]current page.north west) {{\fontsize{12bp}{12bp}\selectfont 的事。”}};
\node[anchor=base west,inner sep=0pt,outer sep=0pt] at ([xshift=72bp,yshift=-234bp]current page.north west) {{\fontsize{12bp}{12bp}\selectfont 【8.23.17】 沛公因为樊哙、张良的缘故，得以脱身返回。}};
\node[anchor=base west,inner sep=0pt,outer sep=0pt] at ([xshift=72bp,yshift=-249bp]current page.north west) {{\fontsize{12bp}{12bp}\selectfont 【8.23.18】 回来后，立刻杀了曹无伤。}};
\node[anchor=base west,inner sep=0pt,outer sep=0pt] at ([xshift=72bp,yshift=-264bp]current page.north west) {{\fontsize{12bp}{12bp}\selectfont 【8.24.1】 项羽向西进军，屠杀无辜，焚毁咸阳秦宫室，所过之处，无不遭到摧残破坏。}};
\node[anchor=base west,inner sep=0pt,outer sep=0pt] at ([xshift=72bp,yshift=-279bp]current page.north west) {{\fontsize{12bp}{12bp}\selectfont 【8.24.2】 秦地的百姓大失所望，然而心里恐惧，不敢不服从。}};
\node[anchor=base west,inner sep=0pt,outer sep=0pt] at ([xshift=72bp,yshift=-294bp]current page.north west) {{\fontsize{12bp}{12bp}\selectfont 【8.25.1】 项羽派人回去报告楚怀王。}};
\node[anchor=base west,inner sep=0pt,outer sep=0pt] at ([xshift=72bp,yshift=-309bp]current page.north west) {{\fontsize{12bp}{12bp}\selectfont 【8.25.2】 楚怀王说：“按照原来的约定办。”}};
\node[anchor=base west,inner sep=0pt,outer sep=0pt] at ([xshift=72bp,yshift=-324bp]current page.north west) {{\fontsize{12bp}{12bp}\selectfont 【8.25.3】 项羽怨恨楚怀王不肯让他与沛公一起西进入关，而派他北上救赵，}};
\node[anchor=base west,inner sep=0pt,outer sep=0pt] at ([xshift=72bp,yshift=-339bp]current page.north west) {{\fontsize{12bp}{12bp}\selectfont 【8.25.4】 在天下诸侯争夺称王关中的约定中落在后面，他就说：“怀王这个人，我家项}};
\node[anchor=base west,inner sep=0pt,outer sep=0pt] at ([xshift=72bp,yshift=-354bp]current page.north west) {{\fontsize{12bp}{12bp}\selectfont 梁所立，没有什么功劳，凭什么主持约定。本来安定天下的，是诸位将领和我项籍。”}};
\node[anchor=base west,inner sep=0pt,outer sep=0pt] at ([xshift=72bp,yshift=-369bp]current page.north west) {{\fontsize{12bp}{12bp}\selectfont 【8.25.5】 就假意推尊楚怀王为义帝，实际上不听从他的命令。}};
\node[anchor=base west,inner sep=0pt,outer sep=0pt] at ([xshift=72bp,yshift=-384bp]current page.north west) {{\fontsize{12bp}{12bp}\selectfont 【8.26.1】 正月，项羽自立为西楚霸王，在梁、楚地区的九个郡称王，建都彭城。}};
\node[anchor=base west,inner sep=0pt,outer sep=0pt] at ([xshift=72bp,yshift=-399bp]current page.north west) {{\fontsize{12bp}{12bp}\selectfont 【8.26.2】 背弃原来的约定，改立沛公为汉王，在巴、蜀、汉中称玉，建都南郑。}};
\node[anchor=base west,inner sep=0pt,outer sep=0pt] at ([xshift=72bp,yshift=-414bp]current page.north west) {{\fontsize{12bp}{12bp}\selectfont 【8.26.3】 把关中瓜分为三，封立秦朝的三个将领：章邯为雍王，建都废丘：司马欣为塞}};
\node[anchor=base west,inner sep=0pt,outer sep=0pt] at ([xshift=72bp,yshift=-429bp]current page.north west) {{\fontsize{12bp}{12bp}\selectfont 王，建都栎阳；董翳为翟王，建都高奴。}};
\end{tikzpicture}
\null
\end{document}

